\chapter{Systems of Differential Equations and Beyond}

In the previous chapter, we wrestled with second-order differential equations
like $y'' + py' + qy = 0$, watching a single mass on a spring trade position for
velocity and settle or oscillate according to one equation for one variable. Up
to now, differential equations have described a solitary quantity changing in
time: a population, a temperature, the charge on a capacitor, the angle of a
swinging pendulum. One unknown function, one equation. Most real systems,
however, involve several quantities changing together. Foxes and rabbits eat and are eaten; their fortunes are
locked together, neither curve makes sense without the other. Two masses joined
by springs trade energy back and forth. Three chemical species react, each
feeding on and feeding the next. In every one of these settings we have not one
unknown function but several, and the rate of change of each depends on
\emph{all} of them at once. We need a \emph{system} of differential equations.

Here two main threads come together. On one side is everything we learned about
derivatives, exponentials, and the scalar ODE $y' = ky$ whose solution is
$e^{kt}$. On the other side is linear
algebra -- vectors, matrices, and above all \textbf{eigenvalues and
eigenvectors}. An eigenvector is a special direction that a matrix merely
stretches without rotating; the eigenvalue is that stretch factor. The
astonishing payoff is this: a whole system of coupled linear equations, no matter
how tangled it looks, can be \emph{uncoupled} by choosing the right coordinates,
and in those coordinates each piece is just $y'=\lambda y$ again. The
eigenvectors are the right coordinates. The eigenvalues are the growth rates.
Solving a linear system generalizes the scalar equation $y'=ky$ to matrix form.

We will build the machinery, draw the pictures that make it intuitive -- the
\emph{phase portraits} that turn a system into a geometric vector field --
classify when equilibria attract and when they repel, and then push at
the edge of the linear world to glimpse what happens when systems go nonlinear:
limit cycles, sensitive dependence, chaos. We end by pointing down the roads
that leave this book and head into the rest of mathematics: partial differential
equations, numerical methods, dynamical systems, and the foundational appendix
that undergirds them all.

\chapterline

\section{What a System Is, and Why You Already Have One}
\label{ch39:what}

A \textbf{first-order system} of differential equations is a collection of
equations, one for the derivative of each unknown function, where each
right-hand side may depend on all the unknowns (and possibly on time):
\[
  \begin{aligned}
    x_1'(t) &= f_1(t, x_1, x_2, \dots, x_n),\\
    x_2'(t) &= f_2(t, x_1, x_2, \dots, x_n),\\
            &\;\;\vdots\\
    x_n'(t) &= f_n(t, x_1, x_2, \dots, x_n).
  \end{aligned}
\]
The word \emph{first-order} means only first derivatives appear. The word
\emph{system} means the equations are bound together -- you cannot solve one
without the others, because each $f_i$ depends on every $x_j$.

It is enormously clarifying to bundle the unknowns into a single vector-valued
function
\[
  \vec{x}(t) = \big(x_1(t), x_2(t), \dots, x_n(t)\big),
\]
and the right-hand sides into a single vector field
$\vec{F}(t,\vec x) = (f_1, \dots, f_n)$. Then the entire system collapses to one
clean line that looks exactly like the scalar case:
\[
  \boxed{\;\vec{x}'(t) = \vec{F}\big(t, \vec{x}(t)\big).\;}
\]
Vector notation here expresses a geometric point of view. A single number
moving on a line obeyed $x' = f(x)$. A single \emph{point} moving through
$n$-dimensional space obeys $\vec x' = \vec F(\vec x)$. The state of the whole
world is one point; the differential equation tells that point, at every moment,
which direction and how fast to move.

\begin{intuitionbox}
Picture the unknowns as the coordinates of a single dot living in a space called
\textbf{phase space}. For the fox--rabbit problem that space is a plane: the
horizontal axis is the rabbit population, the vertical axis the fox population.
The current state of the ecosystem is \emph{one dot} in that plane. The system of
ODEs is a rule that, at every dot, plants a little arrow telling the dot where to
go next. The whole job of solving the system is to follow the arrows -- to let
the dot drift -- and watch the curve it traces. That curve is the history of the
ecosystem.
\end{intuitionbox}

\begin{figure}[H]
\centering
\begin{tikzpicture}[scale=0.92,>={Stealth[length=4pt]}]
  % Panel (a): Phase plane orbit
  \begin{scope}[shift={(0,0)}]
    \draw[fill=blue!3,draw=blue!15,rounded corners=4pt] (-0.8,-0.6) rectangle (5.2,5.0);
    \node[font=\small\bfseries,anchor=north] at (2.2,4.85) {(a) Phase Plane: $(R, F)$ Orbit};
    \draw[->,black!70,thick] (0,0) -- (3.8,0) node[right,font=\scriptsize] {$R$ (Prey)};
    \draw[->,black!70,thick] (0,0) -- (0,3.9) node[above,font=\scriptsize] {$F$ (Predator)};
    % Closed orbit (counterclockwise)
    \draw[curveblue,ultra thick] (2.2,2.0) ellipse (1.6cm and 1.2cm);
    \draw[->,curveblue,very thick] (2.2,0.8) -- (2.6,0.8);
    \draw[->,curveblue,very thick] (3.8,2.0) -- (3.8,2.4);
    \draw[->,curveblue,very thick] (2.2,3.2) -- (1.8,3.2);
    \draw[->,curveblue,very thick] (0.6,2.0) -- (0.6,1.6);
    % Equilibrium
    \filldraw[targetgreen] (2.2,2.0) circle (2pt) node[right=3pt,font=\tiny\bfseries] {$(R^*, F^*)$};
    % Four stages annotation
    \node[font=\tiny,align=center,curveblue!80!black] at (2.2,0.4) {Prey boom};
    \node[font=\tiny,align=center,curveblue!80!black] at (3.8,3.5) {Predators boom};
    \node[font=\tiny,align=center,curveblue!80!black] at (1.0,3.5) {Prey crash};
    \node[font=\tiny,align=center,curveblue!80!black] at (0.9,0.7) {Predators\\starve};
  \end{scope}

  % Panel (b): Time series R(t) and F(t)
  \begin{scope}[shift={(6.4,0)}]
    \draw[fill=blue!3,draw=blue!15,rounded corners=4pt] (-0.6,-0.6) rectangle (5.8,5.0);
    \node[font=\small\bfseries,anchor=north] at (2.6,4.85) {(b) Time Series: Out-of-Phase Waves};
    \draw[->,black!70,thick] (0,0) -- (5.0,0) node[right,font=\scriptsize] {$t$};
    \draw[->,black!70,thick] (0,0) -- (0,3.9) node[above,font=\scriptsize] {Pop.};
    % Prey wave (blue)
    \draw[curveblue,very thick,samples=100,domain=0:4.6]
      plot (\x, {2.0 + 1.2*cos(deg(1.6*\x))});
    % Predator wave (red, lagged by pi/2)
    \draw[tangentred,very thick,samples=100,domain=0:4.6]
      plot (\x, {2.0 + 1.0*sin(deg(1.6*\x))});
    % Legends
    \draw[curveblue,very thick] (0.7,3.9) -- (1.2,3.9);
    \node[right=2pt,font=\scriptsize,curveblue] at (1.2,3.9) {$R(t)$ (Prey)};
    \draw[tangentred,very thick] (3.1,3.9) -- (3.6,3.9);
    \node[right=2pt,font=\scriptsize,tangentred] at (3.6,3.9) {$F(t)$ (Predator)};
    \node[font=\scriptsize,align=center] at (2.6,-0.4) {Predator peaks lag prey peaks\\by a quarter cycle};
  \end{scope}
\end{tikzpicture}
\caption{The predator--prey system in phase space and in time.
(a) In phase space, the ecosystem state $(R, F)$ traverses a closed orbit counterclockwise: abundant prey
feed predators, predator overpopulation crashes the prey, food scarcity starves predators, allowing prey to recover.
(b) Plotted against time, both populations oscillate periodically with identical period, with predator peaks
chronically lagging behind prey peaks.}
\label{ch39:fig-predator-prey}
\end{figure}

\subsection{A first concrete example}

Here is a small but real system. Let $x(t)$ and $y(t)$ be two quantities
obeying
\[
  x' = \phantom{-}x + 2y, \qquad y' = 3x + 2y.
\]
Notice the coupling: $x'$ depends on $y$, and $y'$ depends on $x$. You cannot
integrate the first equation by itself because the $2y$ on its right is an
unknown function you have not found yet. The equations are coupled; solving them requires methods that decouple the variables.

\begin{happeningbox}
Notice the difference between writing the system down and solving it. The
differential equations tell you the instantaneous velocity $(x', y')$ at each
point $(x, y)$, whereas a solution gives an explicit formula $(x(t), y(t))$
telling you where the state is at any time $t$ without having to trace the
step-by-step path in between.
\end{happeningbox}

\section{Reducing Higher-Order Equations to First-Order Systems}
\label{ch39:reduction}

The previous chapter studied second-order equations like $y'' + py' + qy = 0$,
and every higher-order ODE can be rewritten as a first-order system.
This is one of the most useful structural facts in all of differential
equations, and it is direct once you see it.

Take the second-order equation
\[
  y'' = -y
\]
-- the equation of simple harmonic motion, whose solutions you know are sines
and cosines. Introduce a \emph{second unknown} equal to the velocity:
\[
  x_1 = y, \qquad x_2 = y'.
\]
Now ask for the derivatives of these two new unknowns. The first is immediate:
$x_1' = y' = x_2$. The second uses the original equation: $x_2' = y'' = -y =
-x_1$. So the single second-order equation becomes the first-order system
\[
  x_1' = x_2, \qquad x_2' = -x_1.
\]
We have not changed the physics one bit; we have only renamed ``position'' and
``velocity'' as two coordinates and written down how each one changes. The
trick generalizes word for word.

\begin{whybox}
Why does this always work? An $n$th-order equation
$y^{(n)} = G(t, y, y', \dots, y^{(n-1)})$ determines the highest derivative from
all the lower ones. So if we name the lower derivatives as new variables
\[
  x_1 = y,\quad x_2 = y',\quad x_3 = y'',\quad \dots,\quad x_n = y^{(n-1)},
\]
then each variable's derivative is simply the \emph{next} variable in the
list ($x_1' = x_2$, $x_2' = x_3$, and so on) until the very last, whose
derivative is the only place the original equation does any work:
$x_n' = G(t, x_1, \dots, x_n)$. The cascade $x_i' = x_{i+1}$ is pure
bookkeeping; the genuine content of the equation lives entirely in the final
slot. That is why second-order, third-order, and millionth-order equations all
reduce to first-order systems: differentiation just shifts you one notch up a
ladder, and the equation tells you what happens when you fall off the top.
\end{whybox}

\begin{aestheticbox}
This reduction is why we can afford to spend the whole chapter on first-order
systems and lose nothing. A numerical solver written for $\vec x' = \vec F(\vec
x)$ can integrate the orbit of a planet (second order), the vibration of a
bridge (fourth order), or the firing of a neuron (a tangle of several equations)
without modification. That is why almost all general theorems and numerical
algorithms in dynamics are written once, for first-order systems $\vec x' = \vec F(\vec x)$.
\end{aestheticbox}

\begin{example}[Damped oscillator as a system]
\label{ch39:ex-damped}
The equation $y'' + 2y' + 5y = 0$ models a mass on a spring with friction. Set
$x_1 = y$, $x_2 = y'$. Then $x_1' = x_2$, and from the equation
$y'' = -2y' - 5y$, so $x_2' = -5x_1 - 2x_2$. In matrix form,
\[
  \begin{pmatrix} x_1 \\ x_2 \end{pmatrix}'
  =
  \begin{pmatrix} 0 & 1 \\ -5 & -2 \end{pmatrix}
  \begin{pmatrix} x_1 \\ x_2 \end{pmatrix}.
\]
The matrix that appears is called the \textbf{companion matrix} of the
equation -- a standard matrix that translates a single higher-order ODE into a
first-order system by shifting variables along the upper off-diagonal and placing
the original equation's coefficients in the bottom row. We will solve systems of
exactly this shape in a moment; the bottom row is literally the coefficients of
the original ODE, sign-flipped.
\end{example}

\begin{walkbox}
\textbf{Pause and think:} Suppose you have the third-order equation $y''' - y = 0$.
If you set $x_1 = y$, $x_2 = y'$, and $x_3 = y''$, what does the $3\times 3$
companion matrix look like? (The first row gives $x_1' = x_2$, so $(0, 1, 0)$;
the second gives $x_2' = x_3$, so $(0, 0, 1)$; and the bottom row reads the ODE
itself, $x_3' = y''' = y = x_1$, giving $(1, 0, 0)$. Each row shifts one notch
to the right until the bottom row loops back!)
\end{walkbox}

\section{The Matrix Form \texorpdfstring{$\vec{u}' = A\vec{u}$}{u' = Au}}
\label{ch39:matrixform}

From here on we focus on \textbf{linear} systems with constant coefficients,
because they are the ones we can solve completely and because everything we learn
about them governs the local behavior of nonlinear systems too (Section
\ref{ch39:nonlinear}). A linear constant-coefficient system in two unknowns
looks like
\[
  \begin{aligned}
    u_1' &= a u_1 + b u_2,\\
    u_2' &= c u_1 + d u_2,
  \end{aligned}
\]
and packs neatly into the single matrix equation
\[
  \boxed{\;\vec{u}' = A\,\vec{u}, \qquad
  \vec u = \begin{pmatrix} u_1 \\ u_2 \end{pmatrix}, \quad
  A = \begin{pmatrix} a & b \\ c & d \end{pmatrix}.\;}
\]
Here $\vec u'$ means ``differentiate each component'': $\vec u' = (u_1', u_2')$.
The whole system is now a single object: a matrix $A$ acting on a moving vector.

Consider the scalar case as an analogy. If $\vec u$ has just one component, $A$ is
just a number $\lambda$, and the equation is $u' = \lambda u$, whose solution we
have known since the chapter on exponentials:
\[
  u(t) = u(0)\,e^{\lambda t}.
\]
The strategy for solving linear systems extends the exponential trial solution: we test
whether the matrix equation admits solutions of the form $\vec u(t) =
e^{\lambda t}\vec v$ for a constant vector $\vec v$ and scalar rate $\lambda$.
Let us test that hypothesis.

\begin{walkbox}
Suppose that $\vec u(t) = e^{\lambda t}\vec v$ where $\vec v$ is a
constant nonzero vector. Differentiate the left side: since $\vec v$ is constant,
$\vec u'(t) = \lambda e^{\lambda t}\vec v$. Now compute the right side:
$A\vec u(t) = A(e^{\lambda t}\vec v) = e^{\lambda t}(A\vec v)$. For these to be
equal for all $t$, cancel the never-zero scalar $e^{\lambda t}$ and you are left
with
\[
  A\vec v = \lambda\vec v.
\]
\emph{That is the eigenvalue equation.} The requirement that $e^{\lambda t}\vec v$ solve
the differential equation is precisely the condition that $\vec v$ be an eigenvector of
$A$ with eigenvalue $\lambda$. The differential equation is converted directly into a problem of linear algebra.
\end{walkbox}

\begin{connectionbox}
This is the hinge between this book and its companion, \emph{Linear Algebra From
the Ground Up}. There, an eigenvector $\vec v$ of $A$ is defined as a special
direction that $A$ merely \emph{stretches} rather than rotates: $A\vec v =
\lambda\vec v$, with the stretch factor $\lambda$ the eigenvalue. That definition
can feel abstract when you first meet it -- why should anyone care about
directions a matrix only stretches? Here is the answer, arriving from the
calculus side: those are precisely the directions in which a linear system of
ODEs behaves like a single scalar exponential. If you have not yet met
eigenvalues, the chapters ``Eigenvalues and Eigenvectors'' and
``Diagonalization'' in the companion volume develop $\det(A-\lambda I)=0$ and the
solving of $(A-\lambda I)\vec v = \vec 0$ from scratch; we will use those tools
freely below and recall just enough to stay self-contained.
\end{connectionbox}

\subsection{Recipe: eigenvalues and eigenvectors}

For a $2\times 2$ matrix
$A=\begin{psmallmatrix} a & b\\ c & d\end{psmallmatrix}$, the eigenvalues are the
roots $\lambda$ of the \textbf{characteristic equation} -- the polynomial equation
that ensures $(A-\lambda I)$ is singular, so it can squash nonzero vectors into
the zero vector:
\[
  \det(A - \lambda I) = 0
  \quad\Longleftrightarrow\quad
  \lambda^2 - (\underbrace{a+d}_{\text{trace }\tau})\,\lambda
  + (\underbrace{ad-bc}_{\det\,\Delta}) = 0.
\]
Here the \textbf{trace} $\tau = a+d$ is the sum of the main diagonal entries, and
the \textbf{determinant} $\Delta = ad-bc$ measures how the matrix scales areas.
Using the quadratic formula, the eigenvalues are
$\lambda = \tfrac{1}{2}\big(\tau \pm \sqrt{\tau^2 - 4\Delta}\big)$. For each
eigenvalue $\lambda$, the eigenvectors are the nonzero solutions $\vec v$ of the
system $(A-\lambda I)\vec v = \vec 0$. Once we have a pair of
eigenvalue--eigenvector couples $(\lambda_1, \vec v_1)$ and $(\lambda_2, \vec
v_2)$, each gives a solution $e^{\lambda_i t}\vec v_i$, and -- because the system
is linear -- \emph{any combination of solutions is again a solution}.

\begin{theorem}[General solution by eigen-decomposition]
\label{ch39:thm-gensol}
Suppose the $n\times n$ matrix $A$ has $n$ linearly independent eigenvectors
$\vec v_1, \dots, \vec v_n$ with eigenvalues $\lambda_1, \dots, \lambda_n$. Then
\emph{every} solution of $\vec u' = A\vec u$ has the form
\[
  \vec u(t) = c_1 e^{\lambda_1 t}\vec v_1
            + c_2 e^{\lambda_2 t}\vec v_2
            + \cdots
            + c_n e^{\lambda_n t}\vec v_n,
\]
with constants $c_1, \dots, c_n$ fixed by the initial condition $\vec u(0)$.
\end{theorem}

\begin{whybox}
Why does the combination sweep up \emph{every} solution, not just some? Two
reasons, both worth holding onto. First, \emph{superposition}: if $\vec p$ and
$\vec q$ each solve $\vec u' = A\vec u$, then so does $c_1\vec p + c_2\vec q$,
because differentiation and the matrix $A$ are both linear --
$(c_1\vec p + c_2\vec q)' = c_1\vec p' + c_2\vec q' = c_1 A\vec p + c_2 A\vec q =
A(c_1\vec p + c_2\vec q)$. So the listed combinations really are solutions.
Second, \emph{counting}: the existence-and-uniqueness theorem for linear systems
says a solution is pinned down completely by its value at one instant, $\vec
u(0)\in\R^n$ -- that is $n$ numbers of freedom, no more. Our family has exactly
$n$ free constants $c_1,\dots,c_n$, and because the $\vec v_i$ are independent,
choosing them can hit \emph{any} starting vector $\vec u(0) = \sum c_i\vec v_i$.
Because the eigenvectors span $\R^n$, this family satisfies arbitrary initial
conditions $\vec u(0)$, and by the uniqueness theorem it contains every solution.
\end{whybox}

\begin{happeningbox}
The eigenvectors do something decisive: they \emph{uncouple} the
system. Write the starting state in the eigenvector basis, $\vec u(0) = \sum c_i
\vec v_i$. In those coordinates the matrix $A$ stops mixing the components --
along each eigen-direction it acts as pure multiplication by $\lambda_i$. So in
the right coordinates, the tangled system $\vec u' = A\vec u$ falls apart into
$n$ \emph{independent} copies of the scalar equation $z_i' = \lambda_i z_i$, each
solved by $z_i(t) = c_i e^{\lambda_i t}$. The coupling we fought in Section
\ref{ch39:what} was an illusion of bad coordinates. Eigenvectors are the good
coordinates.
\end{happeningbox}

\begin{walkbox}
\textbf{Pause and think:} Look at the general solution
$\vec u(t) = c_1 e^{\lambda_1 t}\vec v_1 + c_2 e^{\lambda_2 t}\vec v_2$.
Suppose $\lambda_1 = -1$ and $\lambda_2 = -4$. As $t \to \infty$, both terms decay to
zero. But which direction does the trajectory hug as it glides into the origin?
(The term $e^{-4t}$ decays four times faster than $e^{-t}$. Long after $e^{-4t}$
has vanished into negligible fractions, the slower $e^{-t}$ term is still alive.
So every typical trajectory flattens out and enters the origin tangent to the
slower eigen-direction $\vec v_1$!)
\end{walkbox}

\section{A Fully Worked \texorpdfstring{$2\times2$}{2x2} System}
\label{ch39:worked}

Let us solve the system from Section \ref{ch39:what} completely, step by step.
\[
  \vec u' = A\vec u, \qquad
  A = \begin{pmatrix} 1 & 2 \\ 3 & 2 \end{pmatrix}.
\]

\textbf{Step 1 -- eigenvalues.} The trace is $\tau = 1 + 2 = 3$ and the
determinant is $\Delta = (1)(2) - (2)(3) = 2 - 6 = -4$. The characteristic
equation is
\[
  \lambda^2 - 3\lambda - 4 = 0 \;\Longrightarrow\; (\lambda - 4)(\lambda + 1) = 0,
\]
so $\lambda_1 = 4$ and $\lambda_2 = -1$. Two real eigenvalues of opposite sign --
already this tells us (Section \ref{ch39:phase}) that the origin will be a
\emph{saddle}.

\textbf{Step 2 -- eigenvector for $\lambda_1 = 4$.} Solve $(A - 4I)\vec v = \vec
0$:
\[
  A - 4I = \begin{pmatrix} 1-4 & 2 \\ 3 & 2-4 \end{pmatrix}
         = \begin{pmatrix} -3 & 2 \\ 3 & -2 \end{pmatrix}.
\]
The top row says $-3v_1 + 2v_2 = 0$, i.e.\ $v_2 = \tfrac{3}{2}v_1$. Choosing
$v_1 = 2$ to clear the fraction gives
$\vec v_1 = \begin{psmallmatrix} 2 \\ 3\end{psmallmatrix}$. (The bottom row
$3v_1 - 2v_2 = 0$ is the same equation -- it must be, because $A-4I$ is singular;
that redundancy is the eigenvalue working as designed.)

\textbf{Step 3 -- eigenvector for $\lambda_2 = -1$.} Solve $(A + I)\vec v = \vec
0$:
\[
  A + I = \begin{pmatrix} 2 & 2 \\ 3 & 3 \end{pmatrix},
  \qquad 2v_1 + 2v_2 = 0 \;\Longrightarrow\; v_2 = -v_1.
\]
Choosing $v_1 = 1$ gives
$\vec v_2 = \begin{psmallmatrix} 1 \\ -1\end{psmallmatrix}$.

\textbf{Step 4 -- assemble.} By Theorem \ref{ch39:thm-gensol},
\[
  \vec u(t)
   = c_1\,e^{4t}\begin{pmatrix} 2 \\ 3 \end{pmatrix}
   + c_2\,e^{-t}\begin{pmatrix} 1 \\ -1 \end{pmatrix}.
\]

\textbf{Step 5 -- fit an initial condition.} Suppose
$\vec u(0) = \begin{psmallmatrix} 3 \\ 2\end{psmallmatrix}$. At $t=0$ the
exponentials are $1$, so we need
\[
  c_1\begin{pmatrix} 2 \\ 3 \end{pmatrix}
  + c_2\begin{pmatrix} 1 \\ -1 \end{pmatrix}
  = \begin{pmatrix} 3 \\ 2 \end{pmatrix}
  \;\Longrightarrow\;
  \begin{cases} 2c_1 + c_2 = 3,\\ 3c_1 - c_2 = 2. \end{cases}
\]
Adding the two equations: $5c_1 = 5$, so $c_1 = 1$, and then $c_2 = 3 - 2 = 1$.
The trajectory through $(3,2)$ is
\[
  \vec u(t) = e^{4t}\begin{pmatrix} 2 \\ 3 \end{pmatrix}
            + e^{-t}\begin{pmatrix} 1 \\ -1 \end{pmatrix}.
\]

\begin{intuitionbox}
Read the answer like a story of two tendencies. The $e^{4t}$ piece grows
explosively and points along $\vec v_1 = (2,3)$: this is the \emph{unstable}
direction, the one the system is fleeing toward. The $e^{-t}$ piece decays and
points along $\vec v_2 = (1,-1)$: the \emph{stable} direction, the one the system
is collapsing onto. Almost immediately the decaying term becomes negligible and
the trajectory becomes nearly parallel to $(2,3)$, hurtling outward. The two
eigenvectors are the grain of the wood; every trajectory bends to follow it.
\end{intuitionbox}

\begin{trapbox}
A frequent slip: writing
$\vec u(t) = e^{\lambda_1 t}\vec v_1 + e^{\lambda_2 t}\vec v_2$ \emph{without} the
constants $c_1, c_2$, then being unable to match the initial condition. The
constants are essential; they allow you to adjust the
solution at the correct starting point. With two eigen-solutions you have two
dials, exactly enough to hit any point in the plane. Drop them and you can
describe only one special trajectory out of the infinitely many.
\end{trapbox}

\section{Phase Portraits: Drawing the Flow}
\label{ch39:phase}

Formulas give exact coordinates, but a picture shows the whole family of
solutions at once. The \textbf{phase
portrait} of a planar system is the $(u_1, u_2)$ plane decorated with the
arrows of the vector field $A\vec u$ and a representative sample of solution
curves -- the \emph{trajectories} -- threading through them. Reading a phase
portrait is reading the destiny of every possible initial state at a single
glance.

Here is the saddle we just solved. The two eigen-directions are drawn as bold
lines; trajectories race away along the unstable direction $(2,3)$ and are
funneled in along the stable direction $(1,-1)$, sweeping out hyperbola-like
curves between them.

\begin{figure}[H]
\centering
\begin{tikzpicture}
\begin{axis}[
  width=10cm, height=8cm,
  axis lines=middle,
  xlabel=$u_1$, ylabel=$u_2$,
  xmin=-3.2, xmax=3.2, ymin=-3.2, ymax=3.2,
  xtick=\empty, ytick=\empty,
  clip=true,
]
% eigen-directions
\addplot[structurepurple, very thick, domain=-1.2:1.2, samples=2]
  ({2*x},{3*x});
\addplot[targetgreen, very thick, domain=-2.5:2.5, samples=2]
  ({x},{-x});
\draw[->,structurepurple,very thick] (axis cs:0,0) -- (axis cs:1.4,2.1);
\draw[->,structurepurple,very thick] (axis cs:0,0) -- (axis cs:-1.4,-2.1);
\draw[->,targetgreen,very thick] (axis cs:2.0,-2.0) -- (axis cs:1.0,-1.0);
\draw[->,targetgreen,very thick] (axis cs:-2.0,2.0) -- (axis cs:-1.0,1.0);
% saddle trajectories u = c1 e^{4s}(2,3) + c2 e^{-s}(1,-1)
\addplot[curveblue, thick, smooth, domain=-1:0.45, samples=80]
  ({0.10*exp(4*x)*2 + 1.0*exp(-x)*1}, {0.10*exp(4*x)*3 + 1.0*exp(-x)*(-1)});
\addplot[curveblue, thick, smooth, domain=-1:0.45, samples=80]
  ({0.10*exp(4*x)*2 - 1.0*exp(-x)*1}, {0.10*exp(4*x)*3 - 1.0*exp(-x)*(-1)});
\addplot[curveblue, thick, smooth, domain=-1:0.45, samples=80]
  ({-0.10*exp(4*x)*2 + 1.0*exp(-x)*1}, {-0.10*exp(4*x)*3 + 1.0*exp(-x)*(-1)});
\addplot[curveblue, thick, smooth, domain=-1:0.45, samples=80]
  ({-0.10*exp(4*x)*2 - 1.0*exp(-x)*1}, {-0.10*exp(4*x)*3 - 1.0*exp(-x)*(-1)});
\node[structurepurple, fill=white, inner sep=1.2pt, anchor=south east] at (axis cs:1.8,2.7) {\small unstable};
\node[targetgreen, fill=white, inner sep=1.2pt, anchor=south west] at (axis cs:1.9,-1.7) {\small stable};
\end{axis}
\end{tikzpicture}
\caption{Phase portrait of the saddle $\vec u' = \protect\begin{psmallmatrix} 1 &
2 \\ 3 & 2\protect\end{psmallmatrix}\vec u$. Eigenvalues $4$ and $-1$. Purple: the
unstable eigen-direction $(2,3)$; green: the stable eigen-direction $(1,-1)$.
Trajectories flow out along purple and in along green.}
\label{ch39:fig-saddle}
\end{figure}

\subsection{The classification: trace and determinant decide everything}

The qualitative shape of a planar phase portrait depends only on the two
eigenvalues of $A$, and those depend only on the trace $\tau = a+d$ and
determinant $\Delta = ad-bc$. Here is the full classification.

\begin{itemize}
\item \textbf{Two real eigenvalues, same sign.} If both are negative, every
trajectory decays to the origin: a stable \textbf{node} (a sink), drawn in
Figure \ref{ch39:fig-node}. If both
positive, everything flees outward: an unstable node (a source). The slower
eigen-direction dominates the shape near the origin.

\item \textbf{Two real eigenvalues, opposite signs.} A \textbf{saddle}, as in
Figure \ref{ch39:fig-saddle}. The origin attracts along one eigen-direction and
repels along the other. Always unstable: almost every trajectory eventually
leaves.

\item \textbf{Complex eigenvalues $\lambda = \alpha \pm i\beta$, $\beta\ne0$.}
The solutions involve $e^{\alpha t}(\cos\beta t, \sin\beta t)$ -- rotation
(from the $\cos,\sin$) times scaling (from $e^{\alpha t}$). If $\alpha<0$,
trajectories spiral inward: a stable \textbf{spiral}. If $\alpha>0$, an unstable
spiral. If $\alpha = 0$ exactly, they neither grow nor shrink but circle
forever: a \textbf{center}, a ring of closed orbits.
\end{itemize}

\begin{whybox}
Where do the rotation and the spiral come from? When the eigenvalues are complex,
$\lambda = \alpha \pm i\beta$, the key relationship is Euler's formula,
$e^{i\beta t} = \cos\beta t + i\sin\beta t$, which we proved back in the
exponentials chapter. The solution
$e^{\lambda t}\vec v = e^{\alpha t}e^{i\beta t}\vec v$ therefore factors into a
real \emph{growth} part $e^{\alpha t}$ and an oscillating \emph{rotation} part
$e^{i\beta t}$. The real eigenvalue gave us pure stretching; the imaginary part
of a complex eigenvalue gives us turning. The imaginary axis is, quite literally,
the axis of rotation in phase space. That is why an undamped spring (eigenvalues
$\pm i\omega$, purely imaginary) orbits forever in a center, while a damped
spring (eigenvalues with negative real part) spirals into rest.
\end{whybox}

The whole classification fits on one diagram in the $(\tau, \Delta)$ plane. The
parabola $\tau^2 = 4\Delta$ is the boundary where the eigenvalues switch from
real to complex.

\begin{figure}[H]
\centering
\begin{tikzpicture}
\begin{axis}[
  width=11cm, height=8cm,
  axis lines=middle,
  xlabel=$\tau$ (trace), ylabel=$\Delta$ (det),
  xlabel style={at={(ticklabel* cs:1)}, anchor=north west},
  xmin=-3.5, xmax=3.5, ymin=-2.2, ymax=4,
  xtick=\empty, ytick=\empty,
  clip=true,
]
% the discriminant parabola tau^2 = 4 Delta
\addplot[tangentred, very thick, domain=-3.4:3.4, samples=80] {x^2/4};
\node[align=left, anchor=west] at (axis cs:0.3,2.8)
  {\small spirals \& centers\\ \small (complex $\lambda$)};
\node[font=\scriptsize, align=left, anchor=west] at (axis cs:0.15,1.5) {centers ($\tau=0$)};
\node[align=center, targetgreen, font=\scriptsize] at (axis cs:-2.2,0.4)
  {stable\\ node};
\node[align=center, trapred, font=\scriptsize] at (axis cs:2.2,0.4)
  {unstable\\ node};
\node[align=center, structurepurple] at (axis cs:1.8,-1.2)
  {\small saddles ($\Delta<0$)};
\node[tangentred, fill=white, inner sep=1.2pt, anchor=south east] at (axis cs:2.7,2.0) {\small $\tau^2{=}4\Delta$};
\end{axis}
\end{tikzpicture}
\caption{The trace--determinant plane. Below the horizontal axis ($\Delta<0$):
saddles. Above it and below the red parabola: real eigenvalues, hence nodes
(stable on the left where $\tau<0$, unstable on the right). Above the parabola:
complex eigenvalues, hence spirals (stable for $\tau<0$, centers on the line
$\tau=0$, unstable for $\tau>0$).}
\label{ch39:fig-tracedet}
\end{figure}

\begin{aestheticbox}
Marvel at the economy. Two numbers -- the trace and determinant of a $2\times2$
matrix -- predict the entire qualitative fate of an infinite family of
trajectories. You need not solve anything to know whether a bridge will settle or
shake apart, whether a population will reach equilibrium or oscillate, whether a
control system is stable. You compute $\tau$ and $\Delta$ and read off the
answer from one picture. This is the gift of the eigenvalue viewpoint: it
converts \emph{dynamics} (hard, time-dependent) into \emph{algebra} (two numbers
on a chart).
\end{aestheticbox}

\begin{walkbox}
\textbf{Pause and think:} Suppose you compute the determinant of a $2\times 2$
system matrix and get $\Delta = -6$. Can the origin ever be a stable equilibrium?
Can it be a spiral? (No! The determinant is the product of the eigenvalues,
$\Delta = \lambda_1 \lambda_2$. If $\Delta < 0$, the eigenvalues cannot be
complex conjugates, because $(\alpha+i\beta)(\alpha-i\beta) = \alpha^2+\beta^2 \ge 0$.
So the eigenvalues must be real with opposite signs: one positive, one negative.
Trajectories will always explode outward along the positive eigen-direction,
making the origin an unstable saddle every single time.)
\end{walkbox}

\subsection{A spiral, drawn}

For concreteness, take $A = \begin{psmallmatrix} -0.3 & -1 \\ 1 &
-0.3\end{psmallmatrix}$. Its eigenvalues are $-0.3 \pm i$, complex with negative
real part: a stable spiral. Every trajectory winds inward to the origin.

\begin{figure}[H]
\centering
\begin{tikzpicture}
\begin{axis}[
  width=9.5cm, height=8cm,
  axis lines=middle,
  xlabel=$u_1$, ylabel=$u_2$,
  xmin=-3, xmax=3, ymin=-3, ymax=3,
  xtick=\empty, ytick=\empty,
]
% trajectory: e^{-0.3 t}(2.7 cos t, 2.7 sin t)
\addplot[curveblue, thick, smooth, domain=0:14, samples=180]
  ({2.7*exp(-0.3*x)*cos(deg(x))}, {2.7*exp(-0.3*x)*sin(deg(x))});
\addplot[targetgreen, thick, smooth, domain=0:14, samples=180]
  ({-1.6*exp(-0.3*x)*cos(deg(x))}, {-1.6*exp(-0.3*x)*sin(deg(x))});
\filldraw[curveblue] (axis cs:2.7,0) circle (1.8pt);
\node[curveblue, anchor=south east] at (axis cs:2.55,0.2) {\small start};
\end{axis}
\end{tikzpicture}
\caption{A stable spiral: $A=\protect\begin{psmallmatrix} -0.3 & -1 \\ 1 &
-0.3\protect\end{psmallmatrix}$, eigenvalues $-0.3\pm i$. Rotation comes from the
imaginary part $\pm i$; the inward decay comes from the real part $-0.3$. Two
trajectories shown.}
\label{ch39:fig-spiral}
\end{figure}

\subsection{A stable node, drawn}

For contrast, take the diagonal matrix
$A = \begin{psmallmatrix} -1 & 0 \\ 0 & -3\end{psmallmatrix}$, with eigenvalues
$-1$ and $-3$ and eigenvectors along the two axes. Both eigenvalues are negative,
so every trajectory decays to the origin -- a stable node. Because the $u_2$
component dies three times faster ($e^{-3t}$ versus $e^{-t}$), trajectories flatten
toward the slow axis $u_1$ as they approach the origin, and that slow direction
sets the shape you see near the origin.

\begin{figure}[H]
\centering
\begin{tikzpicture}
\begin{axis}[
  width=9.5cm, height=8cm,
  axis lines=middle,
  xlabel=$u_1$, ylabel=$u_2$,
  xmin=-3, xmax=3, ymin=-3, ymax=3,
  xtick=\empty, ytick=\empty,
  clip=true,
]
% trajectories u = (c1 e^{-t}, c2 e^{-3t}); parametrize by s=t from 0 to 3.5
\addplot[curveblue, thick, smooth, domain=0:3.5, samples=120]
  ({2.8*exp(-x)}, {2.5*exp(-3*x)});
\addplot[curveblue, thick, smooth, domain=0:3.5, samples=120]
  ({-2.8*exp(-x)}, {2.5*exp(-3*x)});
\addplot[curveblue, thick, smooth, domain=0:3.5, samples=120]
  ({2.8*exp(-x)}, {-2.5*exp(-3*x)});
\addplot[curveblue, thick, smooth, domain=0:3.5, samples=120]
  ({-2.8*exp(-x)}, {-2.5*exp(-3*x)});
\addplot[curveblue, thick, smooth, domain=0:3.5, samples=120]
  ({0.7*exp(-x)}, {2.8*exp(-3*x)});
\addplot[curveblue, thick, smooth, domain=0:3.5, samples=120]
  ({-0.7*exp(-x)}, {-2.8*exp(-3*x)});
\end{axis}
\end{tikzpicture}
\caption{A stable node: $A=\protect\begin{psmallmatrix} -1 & 0 \\ 0 &
-3\protect\end{psmallmatrix}$, eigenvalues $-1$ and $-3$. Every trajectory decays
to the origin. The faster-dying $u_2$ direction flattens each curve onto the slow
$u_1$ axis before it arrives, so the node is ``combed'' along $u_1$.}
\label{ch39:fig-node}
\end{figure}

\section{Stability: When Does the System Settle?}
\label{ch39:stability}

A point $\vec u^{*}$ where $\vec F(\vec u^{*}) = \vec 0$ is an
\textbf{equilibrium} -- plant the state dot there and it never moves, because its
velocity is zero. For the linear system $\vec u' = A\vec u$ the origin is always
an equilibrium (and the only one, if $A$ is invertible). A fundamental practical question
in engineering is: if the system is slightly perturbed from equilibrium,
does it return, or does it diverge?

The eigenvalues answer it cleanly. A solution component behaves like
$e^{\lambda t}$, which shrinks if the real part of $\lambda$ is negative and
grows if it is positive. Hence:

\begin{theorem}[Linear stability]
\label{ch39:thm-stability}
For $\vec u' = A\vec u$, the origin is
\begin{itemize}
\item \textbf{asymptotically stable} (every trajectory $\to \vec 0$) if and only
if \emph{every} eigenvalue of $A$ has negative real part;
\item \textbf{unstable} if some eigenvalue has positive real part;
\item \textbf{marginally stable} (trajectories stay bounded but need not decay)
in the borderline case of purely imaginary eigenvalues, e.g.\ a center.
\end{itemize}
\end{theorem}

\begin{intuitionbox}
The principle is direct: \emph{the system is stable exactly when every growth
rate is a decay rate.} One single eigenvalue poking into the right half of the
complex plane -- one direction in which a perturbation grows -- is enough to
doom stability, because a real trajectory generically has a nonzero component
along every eigen-direction, and the growing one eventually swamps the rest.
Stability is unanimity: \emph{all} must decay.
\end{intuitionbox}

\begin{applicationbox}
\textbf{Flight stability and eigenvalue analysis.} An aircraft in
steady flight is an equilibrium of a system of differential equations governing
its pitch, roll, and yaw. Engineers compute the eigenvalues of the linearized
dynamics. If all have negative real part, a gust that tips the plane decays away
and the plane self-corrects -- the famous \emph{phugoid} mode is a slow stable
spiral you can sometimes feel as a gentle porpoising. If an eigenvalue creeps
into the right half-plane, the aircraft is unstable and needs a control system
(itself adding feedback terms to $A$, shifting the eigenvalues leftward) to fly
at all. Modern fighter jets are deliberately built unstable for agility and are
flyable only because a computer moves the eigenvalues thousands of times a
second. Eigenvalue placement is a multi-billion-dollar industry.
\end{applicationbox}

\begin{historybox}
The geometric and stability theory of differential equations was forged in the
late nineteenth century when classical analytical methods proved insufficient for
complex systems. For two hundred years after Newton, mathematicians sought exact
algebraic formulas for every moving system. They solved the two-body problem (a single
planet orbiting the Sun) with clean Keplerian ellipses, but the three-body problem
(the Sun, Earth, and Moon tugging on one another) defeated every attempt. At the same
time, the Industrial Revolution produced high-speed steam engines whose
mechanical regulators (governors) suffered from violent, destructive
oscillations called ``hunting'' -- machines that shook themselves apart because
their feedback loops overcorrected.

Armed only with paper-and-pencil algebra, linear approximations, and
perturbation series that often quietly diverged, mathematicians were forced to
invent new tools. James Clerk Maxwell opened stability theory in 1868 with ``On
Governors,'' analyzing when linear feedback stabilizes an engine. Edward Routh
(1876) and Adolf Hurwitz (1895) devised algebraic recipes to test whether all
roots of a polynomial have negative real parts without calculating them. In 1892,
the Russian mathematician Aleksandr Lyapunov published \emph{The General
Problem of the Stability of Motion}, introducing what we now call Lyapunov
functions: generalized energy-like functions that prove a system settles back to
rest without needing explicit solutions.

A major conceptual development arrived with Henri Poincar\'e. In 1889, King Oscar II
of Sweden sponsored an international prize to determine whether the solar system is stable. Poincar\'e submitted a
memoir claiming a proof and won the gold medal. But while the paper was being
typeset for the journal \emph{Acta Mathematica}, the young editor Edvard
Phragm\'en spotted an unexplained gap. Investigating it, Poincar\'e discovered to
his horror that his proof had collapsed: trajectories did not glide smoothly
along regular surfaces but crossed each other in an infinite, self-intersecting
tangle. Poincar\'e paid thousands of francs out of his own pocket to buy back
and destroy the printed copies of the journal, rewrote the memoir, and in doing
so discovered what we now call deterministic chaos.

Poincar\'e abandoned the hunt for closed-form formulas and founded the modern
qualitative theory of dynamical systems: plotting trajectories directly in phase
space, classifying equilibria, and tracking geometric flows. Yet the field
remains alive with unsolved mysteries. The second half of Hilbert's 16th
Problem, posed in 1900, asks for the maximum number and possible relative
arrangements of limit cycles (isolated periodic orbits) for polynomial systems
$x'=P(x,y)$, $y'=Q(x,y)$. Even for quadratic polynomials (degree two), it
remains completely unsolved today.
\end{historybox}

\section{A Glimpse Beyond the Linear World}
\label{ch39:nonlinear}

Everything so far rested on linearity -- the right-hand side was $A\vec u$, a
matrix times the state. Real systems are usually \emph{nonlinear}: the
right-hand side bends. Foxes eating rabbits, the swing of a pendulum at large
angles, the weather. We cannot solve these in closed form. But two ideas carry an
astonishing distance.

\subsection{Linearization: zoom in and the curve looks straight}

Near an equilibrium $\vec u^{*}$, expand the vector field in a Taylor series (the
multivariable kind from the partial-derivatives chapters). The constant term
vanishes because $\vec F(\vec u^{*})=\vec 0$, and to first order
\[
  \vec F(\vec u^{*} + \vec h) \approx J\,\vec h,
  \qquad
  J = \begin{pmatrix}
    \pdv{F_1}{u_1} & \pdv{F_1}{u_2} \\[4pt]
    \pdv{F_2}{u_1} & \pdv{F_2}{u_2}
  \end{pmatrix}_{\vec u^{*}},
\]
where $J$ is the \textbf{Jacobian matrix} -- a grid of all first partial
derivatives that acts as the multivariable best linear approximation, the
higher-dimensional version of the single derivative $f'(x)$. So \emph{near an
equilibrium, a nonlinear system looks like the linear system $\vec h' = J\vec h$},
and the eigenvalues of $J$ classify the equilibrium just as before.

\begin{whybox}
This is the deepest reason the linear theory was worth all our effort. A
nonlinear vector field, viewed under a strong enough microscope centered on an
equilibrium, is indistinguishable from its linearization -- exactly as a smooth
curve, magnified near a point, becomes its tangent line. The Hartman--Grobman
theorem makes this precise: provided no eigenvalue of $J$ sits on the imaginary
axis, the nonlinear phase portrait near the equilibrium is a gently warped copy
of the linear one. Saddles stay saddles, sinks stay sinks, spirals stay spirals.
The whole trace--determinant classification governs nonlinear systems too -- locally.
Calculus's founding move, ``replace the curve by its tangent,'' returns one last
time, now in the shape of replacing a vector field by its Jacobian.
\end{whybox}

\begin{example}[The pendulum]
\label{ch39:ex-pendulum}
The undamped pendulum $\theta'' = -\sin\theta$ becomes the system $x_1' = x_2$,
$x_2' = -\sin x_1$ (with $x_1=\theta$ the angle and $x_2 = \theta'$ the angular
velocity). Its Jacobian is $J = \begin{psmallmatrix} 0 & 1 \\ -\cos x_1 &
0\end{psmallmatrix}$. At the bottom-hanging rest state $x_1 = 0$:
$J = \begin{psmallmatrix} 0 & 1 \\ -1 & 0\end{psmallmatrix}$, eigenvalues
$\pm i$ -- a \emph{center}, the pendulum swinging back and forth forever without
friction. At the inverted balance point $x_1 = \pi$:
$J = \begin{psmallmatrix} 0 & 1 \\ 1 & 0\end{psmallmatrix}$, eigenvalues
$\pm 1$ -- a \emph{saddle}, unstable, exactly matching the physical fact that
balancing a pendulum upright is hopeless: the tiniest nudge runs away along the
unstable eigen-direction.

Linearization tells us what happens near each equilibrium, but look at Figure
\ref{ch39:fig-pendulum} to see how the whole space weaves together. The motion
is governed by total mechanical energy
$E = \tfrac{1}{2}(\theta')^2 - \cos\theta$. At low energy ($E < 1$), the pendulum
is trapped in the potential well, tracing closed blue loops around the center
$(0,0)$ in periodic oscillations. At high energy ($E > 1$), the pendulum has
enough momentum to crest the top, whirling endlessly in wavy purple tracks that
never reverse direction.

Dividing these two behaviors is a razor-thin boundary curve called the
\textbf{separatrix} (from the Latin for ``divider''), shown in red. The
separatrix runs directly between the inverted saddles $(\pm\pi, 0)$ with energy
$E=1$, obeying $\theta' = \pm 2\cos(\theta/2)$. A trajectory launched precisely
on the separatrix would coast upward, slowing exponentially and taking an
infinite amount of time to come to rest balanced vertically on top.
\end{example}

\begin{figure}[H]
\centering
\begin{tikzpicture}
\begin{axis}[
  width=11.5cm, height=7.5cm,
  axis lines=middle,
  xlabel=$\theta$, ylabel=$\theta'$,
  xmin=-4.2, xmax=4.2, ymin=-2.7, ymax=2.7,
  xtick={-3.1415, 3.1415},
  xticklabels={$-\pi$, $\pi$},
  ytick=\empty,
  clip mode=individual,
  clip=true,
]
% Whirling trajectories (E > 1, passing over the top)
\addplot[structurepurple, thick, smooth, domain=-4.2:4.2, samples=100]
  {sqrt(2*(1.4 + cos(deg(x))))};
\addplot[structurepurple, thick, smooth, domain=-4.2:4.2, samples=100]
  {-sqrt(2*(1.4 + cos(deg(x))))};
% Separatrices (E = 1, connecting the saddles at +/- pi)
\addplot[tangentred, very thick, smooth, domain=-3.1415:3.1415, samples=100]
  {2*cos(deg(x/2))};
\addplot[tangentred, very thick, smooth, domain=-3.1415:3.1415, samples=100]
  {-2*cos(deg(x/2))};
% Oscillating orbits (E < 1, swinging back and forth)
\addplot[curveblue, thick, smooth, domain=-1.57:1.57, samples=80]
  {sqrt(2*(cos(deg(x))))};
\addplot[curveblue, thick, smooth, domain=-1.57:1.57, samples=80]
  {-sqrt(2*(cos(deg(x))))};
\addplot[curveblue, thick, smooth, domain=-0.9:0.9, samples=60]
  {sqrt(2*(cos(deg(x)) - cos(deg(0.9))))};
\addplot[curveblue, thick, smooth, domain=-0.9:0.9, samples=60]
  {-sqrt(2*(cos(deg(x)) - cos(deg(0.9))))};
% Equilibria points
\filldraw[targetgreen] (axis cs:0,0) circle (2.2pt);
\filldraw[trapred] (axis cs:3.1415,0) circle (2.2pt);
\filldraw[trapred] (axis cs:-3.1415,0) circle (2.2pt);
% Labels
\node[targetgreen, font=\scriptsize, fill=white, inner sep=1.2pt, anchor=north east] at (axis cs:-0.15,-0.12) {center $(0,0)$};
\node[trapred, font=\scriptsize, fill=white, inner sep=1.2pt, anchor=south east] at (axis cs:3.1415,0.15) {saddle $(\pi,0)$};
\node[trapred, font=\scriptsize, fill=white, inner sep=1.2pt, anchor=south west] at (axis cs:-3.1415,0.15) {saddle $(-\pi,0)$};
\node[tangentred, font=\scriptsize, fill=white, inner sep=1pt, anchor=south east] at (axis cs:-0.4, 2.05) {separatrix ($E=1$)};
\node[curveblue, font=\scriptsize, fill=white, inner sep=1pt] at (axis cs:0, 0.46) {swinging};
\node[structurepurple, font=\scriptsize, fill=white, inner sep=1pt, anchor=south west] at (axis cs:1.8, 2.1) {whirling};
\end{axis}
\end{tikzpicture}
\caption{Phase portrait of the nonlinear pendulum $\theta'' = -\sin\theta$.
Green dot: the stable hanging rest state $(0,0)$, surrounded by closed blue orbits (periodic swinging).
Red dots: unstable inverted balance points $(\pm\pi, 0)$ acting as saddles.
The red separatrix divides swinging back and forth from continuous purple whirling over the top.}
\label{ch39:fig-pendulum}
\end{figure}

\begin{walkbox}
\textbf{Pause and think:} At the bottom rest point $\theta = 0$, the undamped
pendulum has eigenvalues $\pm i$ (a center). In a real room with friction, how
do the eigenvalues change, and what happens to the phase portrait? (Air
resistance introduces a velocity drag term $-b\theta'$, which adds $-b$ to the
trace. The eigenvalues become $-b/2 \pm i\sqrt{1 - b^2/4}$, shifting their real
part into negative territory. The concentric closed orbits become inward spirals:
every swinging pendulum eventually winds down and comes to rest at the bottom!)
\end{walkbox}

\subsection{What linearization cannot see: limit cycles and chaos}

Linearization approximates a nonlinear system locally by its tangent dynamics. Globally,
nonlinear systems exhibit behaviors that do not appear in linear systems.

A \textbf{limit cycle} is an isolated closed loop that nearby trajectories spiral
\emph{onto} -- a self-sustaining oscillation with its own preferred amplitude,
like a heartbeat or the buzz of a violin string. The Van der Pol oscillator
$x'' - \mu(1-x^2)x' + x = 0$ is the classic example; no combination of
eigenvectors can produce it, because a linear center has a whole continuum of
orbits with no preferred size, while a limit cycle picks out exactly one.

And in three or more dimensions, nonlinear systems can be \textbf{chaotic}:
trajectories that never repeat, never settle, and depend so sensitively on their
starting point that two neighboring states diverge exponentially. The most famous
is the Lorenz system,
\[
  x' = \sigma(y-x), \qquad y' = x(\rho - z) - y, \qquad z' = xy - \beta z,
\]
which Edward Lorenz extracted in 1963 from a simplified model of atmospheric
convection. Its trajectories are forever drawn toward a geometric object called
an \textbf{attractor} -- a shape in phase space that pulls in all nearby
trajectories and traps them, yet on which they wander forever without repeating
or crossing. This is the celebrated \emph{Lorenz attractor}, looping around two
swirling lobes in an order no formula predicts.

\begin{applicationbox}
\textbf{The butterfly effect and the limits of forecasting.} Lorenz discovered
chaos by accident. Rerunning a weather simulation on his 1960s computer, he
typed in a midstream value rounded to three decimals ($0.506$) instead of the
full six the machine stored ($0.506127$), expecting the exact same forecast.
Within simulated weeks the two runs had diverged completely, with nothing in
common. That tiny rounding error, amplified by the system's exponential
sensitivity, rewrote the future. Lorenz named the lesson the \textbf{butterfly
effect}: a butterfly flapping its wings in Brazil might set off a tornado in
Texas -- not because the flap carries immense energy, but because the atmosphere
continually magnifies microscopic differences. Long-range weather prediction is
therefore severely constrained in practice, no matter how powerful the computer or how
dense the weather sensors.
\end{applicationbox}

\begin{aestheticbox}
The Lorenz system separates \emph{determinism} from \emph{predictability}. Nothing in the
equations is random: each initial point determines a single, unique trajectory for all time.
Yet because nearby trajectories separate exponentially fast, any tiny measurement uncertainty
in the starting state grows until it covers the whole attractor. Exact laws do not guarantee
long-range forecasts when initial measurements have finite precision.
\end{aestheticbox}

\section{Where to Go Next}
\label{ch39:next}

This chapter sits at the frontier of what a first calculus book can cover. Three
roads lead onward, and you are now equipped to walk all of them.

\textbf{Partial differential equations.} Every system here had a single
independent variable, time, so the unknowns were functions of $t$ alone --
\emph{ordinary} differential equations. But heat spreads through a \emph{rod},
varying in both position and time; a drumhead vibrates across a \emph{surface};
a quantum wavefunction fills \emph{space}. These obey \textbf{partial}
differential equations, where the unknown depends on several variables and
\emph{partial} derivatives appear -- the heat equation
$\pdv{u}{t} = k\,\frac{\partial^2 u}{\partial x^2}$, the wave equation, Laplace's
equation, the Schrodinger equation. The partial derivatives of the multivariable
chapters and the series of the Taylor chapters (Fourier series, their cousin) are
exactly the tools that crack them open.

\textbf{Numerical methods.} Almost no differential equation arising in practice
can be solved by a formula. We solve them by \emph{marching}: from the state now,
the equation gives the velocity now, so we step a tiny $\Delta t$ forward and
repeat. The crudest version is Euler's method,
$\vec u_{n+1} = \vec u_n + \Delta t\,\vec F(\vec u_n)$; the standard tool is the
fourth-order Runge--Kutta method, which samples the slope four times per step to
cancel error. Every weather forecast, climate model, spacecraft trajectory, and
animated-film cloth simulation is a differential equation being marched forward by
a computer.

\textbf{Dynamical systems.} The phase-portrait, stability, and chaos ideas we
touched are the doorway to a vast modern field -- bifurcations (where an
equilibrium suddenly changes character as a parameter crosses a threshold),
strange attractors, ergodic theory, and the geometry of flows.

\begin{connectionbox}
Notice how completely linear algebra and calculus have fused by the end of this
book. We began with the slope of one curve; we end with the eigenvalues of a
matrix dictating the fate of a whole flowing space. The two subjects are not
neighbors but the same subject seen from two sides: calculus is linear algebra
done to functions and rates, linear algebra is the skeleton calculus hangs on.
If this chapter sent you hunting through the companion \emph{Linear Algebra From
the Ground Up} for eigenvalues, diagonalization, and the matrix exponential
$e^{At}$ (which packages \emph{all} solutions of $\vec u'=A\vec u$ into a single
object $\vec u(t) = e^{At}\vec u(0)$), then the two books have done together what
neither could alone.
\end{connectionbox}

\begin{recoverybox}
Looking back at the progression: we expressed coupled equations in matrix form $\vec u' = A\vec u$; we tested the solution form $e^{\lambda t}\vec v$, leading directly to the eigenvalue problem $A\vec v = \lambda \vec v$; we found that eigenvectors decouple the system into scalar components; we drew the flow as a phase portrait and classified it by trace and determinant; we determined stability from the signs of eigenvalue real parts; and we observed how nonlinear systems can be analyzed locally via their Jacobians before displaying global phenomena like limit cycles and chaos.
\end{recoverybox}

\textbf{Where to go next in this volume.}
Two appendices close the book. The first reviews the algebraic identities, inequalities, absolute-value geometry, functions, and proof techniques (direct proof, contradiction, and induction) used throughout the text. The second develops the completeness of the real numbers -- the least-upper-bound axiom that underpins every existence theorem in calculus.

\section{Exercises}
\label{ch39:exercises}

\begin{exercise}[Reduce to a system]
Convert each higher-order equation to a first-order system $\vec u' = A\vec u$
(or $\vec x' = \vec F(\vec x)$ if nonlinear), naming your new variables
explicitly. (a) $y'' - 3y' + 2y = 0$. (b) $y''' = y$. (c) $y'' + \sin y = 0$.
For (a), write down the companion matrix and state its trace and determinant.
\end{exercise}

\begin{exercise}[Eigen-solve a system]
For $A = \begin{psmallmatrix} 2 & 1 \\ 1 & 2\end{psmallmatrix}$, find the
eigenvalues and eigenvectors, write the general solution of $\vec u' = A\vec u$,
and find the particular solution with $\vec u(0) = \begin{psmallmatrix} 1 \\
3\end{psmallmatrix}$. Classify the origin.
\end{exercise}

\begin{exercise}[A saddle]
The matrix $A = \begin{psmallmatrix} 1 & 3 \\ 3 & 1\end{psmallmatrix}$ has
eigenvalues $4$ and $-2$. Find the eigenvectors and sketch the phase portrait by
hand, marking the stable and unstable directions and the flow along each.
\end{exercise}

\begin{exercise}[Complex eigenvalues]
For $A = \begin{psmallmatrix} 0 & -2 \\ 2 & 0\end{psmallmatrix}$, show the
eigenvalues are $\pm 2i$. What kind of equilibrium is the origin? Without solving,
describe the shape of every trajectory. Then verify that $u_1 = \cos 2t$, $u_2 =
\sin 2t$ is a solution.
\end{exercise}

\begin{exercise}[Trace--determinant reading]
For each matrix, compute $\tau$ and $\Delta$ and classify the origin
(stable/unstable node, saddle, stable/unstable spiral, or center) \emph{without}
finding the eigenvectors. (a) $\begin{psmallmatrix} -1 & 0 \\ 0 &
-3\end{psmallmatrix}$. (b) $\begin{psmallmatrix} -1 & 4 \\ -1 &
-1\end{psmallmatrix}$. (c) $\begin{psmallmatrix} 2 & 1 \\ 0 &
3\end{psmallmatrix}$. (d) $\begin{psmallmatrix} 0 & 1 \\ -4 &
0\end{psmallmatrix}$.
\end{exercise}

\begin{exercise}[Stability hinges on a parameter]
Consider $A = \begin{psmallmatrix} -1 & 1 \\ -1 & k\end{psmallmatrix}$ with $k$
a real parameter. Compute the trace and determinant, and find the range of $k$
for which the origin is asymptotically stable. What happens at the value of $k$
where stability is lost?
\end{exercise}

\begin{exercise}[Damped oscillator, both ways]
Return to $y'' + 2y' + 5y = 0$ from Example \ref{ch39:ex-damped}. (a) Find the
eigenvalues of the companion matrix. (b) Confirm they are complex with negative
real part, and say what kind of equilibrium and what physical behavior this
predicts. (c) Relate the real and imaginary parts of the eigenvalues to the decay
rate and oscillation frequency of the solution.
\end{exercise}

\begin{exercise}[Linearize a nonlinear system]
The system $x' = x(3 - x - y)$, $y' = y(2 - x)$ models two competing species.
(a) Find all equilibria (there are three, including one in the interior with both
populations positive). (b) Compute the Jacobian $J$ at the interior equilibrium.
(c) Find its eigenvalues and classify that equilibrium. Is coexistence stable?
\end{exercise}

\begin{exercise}[Pendulum equilibria]
Using Example \ref{ch39:ex-pendulum}, add light damping: $x_1' = x_2$, $x_2' =
-\sin x_1 - 0.2\,x_2$. Recompute the Jacobian at the hanging rest state $(0,0)$.
How do its eigenvalues differ from the undamped case, and what does the change
predict about the pendulum's motion?
\end{exercise}

\begin{exercise}[Conceptual: why eigenvectors?]
In your own words, explain why guessing a solution of the form $\vec u(t) =
e^{\lambda t}\vec v$ forces $\vec v$ to be an eigenvector of $A$. Then explain
why a $2\times2$ system with two independent eigenvectors has \emph{exactly} the
right number of free constants to match any initial condition.
\end{exercise}

\begin{exercise}[A taste of chaos, by hand]
You cannot solve the Lorenz system, but you can march it. With $\sigma=10$,
$\rho=28$, $\beta=8/3$, and step $\Delta t = 0.01$, do \emph{three} Euler steps
starting from $(x,y,z) = (1,1,1)$, reporting all coordinates each step. Then
repeat from $(1.001, 1, 1)$. Comment on how the two runs compare after three
steps -- and explain why, after a few thousand steps, they would have nothing in
common.
\end{exercise}
